% ===== FILE 4/9: Lan truyền gradient, Tích luỹ gradient, Triệt tiêu gradient, ADC tổng quan, Clone =====

% --- Slide: Lan truyền gradient ---
\begin{frame}{Lan truyền gradient: từ loss ngược về từng tham số Gaussian}
\scriptsize
Toàn bộ pipeline vừa trình bày --- chiếu, EWA, alpha blending --- là một chuỗi phép toán \textbf{khả vi} (differentiable),
nên PyTorch có thể tự động tính đạo hàm của loss theo bất kỳ tham số nào bằng quy tắc chuỗi (chain rule):
\[
\frac{\partial\mathcal{L}}{\partial\theta}=\sum_{x}\frac{\partial\mathcal{L}}{\partial C(x)}\cdot\frac{\partial C(x)}{\partial\theta},
\qquad \theta\in\{\mu, s, q, \alpha, k_{\ell m}\}
\]
\begin{itemize}\itemsep1pt
    \item Với mỗi pixel $x$ mà một Gaussian có góp mặt, đạo hàm $\partial C(x)/\partial\theta$ lan ngược qua đúng chuỗi
    các bước đã tính thuận: alpha blending $\to \alpha_n(x)=\min(0.99,\alpha_nG_n(x)) \to G_n(x) \to \Sigma',\mu'$ (EWA)
    $\to \Sigma,\mu$ (không gian 3D) $\to s,q$ (tham số hoá hiệp phương sai).
    \item \textbf{Quan sát then chốt cho phần sau:} gradient của \textbf{vị trí chiếu} $\partial\mathcal{L}/\partial\mu'$
    (trong không gian màn hình) có xu hướng \textbf{tập trung mạnh ở vùng biên} của Gaussian trong ảnh --- nơi cường độ
    ảnh thay đổi nhanh nhất --- chứ không rải đều khắp cả footprint. Điều này hợp lý về trực giác: ở vùng phẳng, sai một
    chút vị trí Gaussian không ảnh hưởng nhiều tới màu render; ở biên/cạnh sắc, một sai lệch nhỏ về vị trí gây thay đổi
    màu đáng kể ngay lập tức.
    \item Chính đại lượng $\partial\mathcal{L}/\partial\mu'$ này là \textbf{tín hiệu} mà Adaptive Density Control dùng để
    quyết định "Gaussian này có đang thiếu/thừa hình học không" (2 slide sau) --- nhưng bản thân tín hiệu này cũng ẩn
    chứa một vấn đề tinh vi, sẽ nói ngay sau đây.
\end{itemize}
\begin{center}
\includegraphics[width=0.36\linewidth]{ch6_edge_gradient.png}
\end{center}
\end{frame}

% --- Slide: Tích luỹ gradient view-space ---
\begin{frame}{Tích luỹ gradient qua nhiều view: lấy CHUẨN trước, cộng sau}
\scriptsize
Một Gaussian thường được camera train nhìn thấy từ \textbf{nhiều góc khác nhau} trong một epoch (một lượt qua hết mọi
camera). Để quyết định densify, cần tổng hợp gradient thu được từ tất cả các lần render đó. Hàm
\texttt{add\_densification\_stats} (\texttt{scene/gaussian\_model.py:1054--1057}) làm điều này như sau:
\[
\boxed{\;\texttt{xyz\_gradient\_accum}_i \mathrel{+}= \big\lVert \nabla_{\mu'_i}\mathcal{L}\big\rVert_{2}\;},\qquad
\bar g_i=\frac{\texttt{xyz\_gradient\_accum}_i}{\texttt{denom}_i}
\]
với \texttt{denom}$_i$ là số lần Gaussian $i$ được nhìn thấy (và cộng dồn) tính tới thời điểm đó.
\begin{itemize}\itemsep1pt
    \item \textbf{Điểm mấu chốt cần nhớ:} phép \textbf{lấy chuẩn} $\lVert\cdot\rVert$ diễn ra \textbf{ngay tại mỗi lần
    render} (mỗi view), \textbf{trước khi} cộng dồn qua \texttt{denom} --- không phải cộng các \emph{vector} gradient
    của nhiều view lại với nhau rồi mới lấy chuẩn một lần ở cuối.
    \item Hệ quả trực tiếp: vì mỗi số hạng được cộng vào \texttt{xyz\_gradient\_accum} luôn là một \textbf{số không âm}
    (chuẩn của một vector luôn $\ge0$), gradient từ các view \textbf{khác nhau tuyệt đối không thể triệt tiêu lẫn nhau}
    trong phép cộng này.
    \item Repo này còn tích luỹ song song một kênh \textbf{trị tuyệt đối} riêng biệt, \texttt{xyz\_gradient\_accum\_abs}
    --- dùng làm tín hiệu cho ngưỡng \texttt{grad\_abs\_thresh} khi quyết định \textbf{split}, tách biệt hoàn toàn với
    ngưỡng \texttt{grad\_thresh} dùng cho \textbf{clone}. Vì sao cần tới hai kênh gradient khác nhau --- xem slide kế tiếp.
\end{itemize}
\begin{center}
\includegraphics[width=0.36\linewidth]{fig_02_accum_denom.png}
\end{center}
\end{frame}

% --- Slide: Triệt tiêu gradient ---
\begin{frame}{Vấn đề "triệt tiêu gradient": xảy ra GIỮA CÁC PIXEL, không phải giữa các view}
\scriptsize
\textbf{Cách hiểu sai thường gặp:} "gradient từ các camera khác nhau kéo ngược chiều nhau nên triệt tiêu lẫn nhau".
Cách hiểu này \textbf{sai} với chính pipeline đã trình bày --- vì chuẩn $\lVert\cdot\rVert$ đã được lấy \textbf{trước
khi} cộng qua các view (slide trước), phép cộng \emph{giữa các view} không thể nào cho ra một tổng nhỏ hơn từng số hạng.
\begin{itemize}\itemsep1pt
    \item \textbf{Sự thật:} hiện tượng triệt tiêu thật sự xảy ra \emph{bên trong một lần render duy nhất}, \textbf{giữa
    các pixel} thuộc cùng footprint của \textbf{một} Gaussian. Ví dụ: nếu ảnh thật có một cạnh chạy ngang qua giữa
    Gaussian, nửa pixel bên trái có thể "kéo" $\mu'$ sang trái để khớp tốt hơn, nửa bên phải lại "kéo" sang phải ---
    khi \textbf{cộng vector gradient của từng pixel lại} (trước khi lấy chuẩn, ngay trong một lần render), hai lực kéo
    ngược chiều này có thể gần như \textbf{triệt tiêu nhau}, cho ra một vector tổng rất nhỏ dù bản thân mỗi pixel đều có
    gradient đáng kể.
    \item \textbf{Hệ quả nguy hiểm:} một Gaussian đang thực sự "dưới tái tạo" (thiếu chi tiết hình học ở đó) nhưng đối
    xứng theo một trục nào đó có thể bị đánh giá nhầm là "gradient thấp $\Rightarrow$ không cần densify" --- một điểm mù
    của cơ chế chỉ dùng gradient có dấu.
    \item \textbf{Cách khắc phục (kỹ thuật AbsGS):} renderer CUDA tính thêm, ngay trong backward pass, một kênh gradient
    \textbf{trị tuyệt đối theo từng pixel} \emph{trước khi} các dấu đối lập có cơ hội triệt tiêu nhau --- chính là kênh
    \texttt{xyz\_gradient\_accum\_abs} đã nhắc ở slide trước.
\end{itemize}
\begin{center}
\includegraphics[width=0.36\linewidth]{fig_03_grad_cancel.png}
\end{center}
\end{frame}

% --- Slide: ADC tổng quan ---
\begin{frame}{Adaptive Density Control (ADC): khung chung}
\scriptsize
Có gradient tích luỹ $\bar g$ (và $\bar g_{\text{abs}}$) rồi, cứ mỗi \texttt{densification\_interval}$=100$ iteration
(bắt đầu từ \texttt{densify\_from\_iter}$=500$, tức lần đầu ở iteration $600$, tới \texttt{densify\_until\_iter}$=15\,000$),
3DGS thực hiện $3$ loại thao tác:
\begin{center}
\begin{tabular}{lll}
\toprule
\textbf{Thao tác} & \textbf{Điều kiện kích hoạt (ý tưởng)} & \textbf{Mục đích}\\
\midrule
\textbf{Clone} & gradient cao, Gaussian đang \emph{nhỏ} & bù \emph{under-reconstruction} (thiếu hình học)\\
\textbf{Split} & gradient cao, Gaussian đang \emph{to} & bù \emph{over-reconstruction} (Gaussian quá to, làm mờ chi tiết)\\
\textbf{Prune} & opacity quá thấp, hoặc quá to & dọn Gaussian không còn đóng góp\\
\bottomrule
\end{tabular}
\end{center}
"To/nhỏ" được so với ngưỡng \texttt{percent\_dense}$\times$\texttt{extent} (scene extent đã định nghĩa ở phần đầu):
Gaussian có trục dài nhất $\le$ ngưỡng này thì thuộc diện \textbf{clone}, lớn hơn thì thuộc diện \textbf{split}.
\begin{itemize}\itemsep1pt
    \item Trực giác: một Gaussian \textbf{nhỏ} nhưng vẫn còn gradient cao có nghĩa là "còn thiếu chi tiết ở đây, cần
    thêm quân số" $\to$ nhân bản. Một Gaussian \textbf{to} mà vẫn còn gradient cao nghĩa là "đang cố khớp một vùng có
    chi tiết mà một Gaussian không đủ" $\to$ chia nhỏ ra.
\end{itemize}
\end{frame}

% --- Slide: Clone ---
\begin{frame}{Clone: sao chép y nguyên, KHÔNG lệch vị trí}
\scriptsize
Khi một Gaussian nhỏ có gradient vượt ngưỡng, thao tác \texttt{densify\_and\_clone} thực hiện điều tưởng chừng đơn giản
đến bất ngờ: \textbf{sao chép y hệt} mọi tham số sang một Gaussian con mới --- \textbf{cùng vị trí}, cùng scale, cùng
rotation, cùng màu:
\[
\theta_{\text{con}} = \theta_{\text{cha}} \qquad (\text{không có bất kỳ }+\varepsilon\text{ hay lấy mẫu ngẫu nhiên nào})
\]
\begin{itemize}\itemsep1pt
    \item \textbf{Câu hỏi tự nhiên:} nếu hai bản sao giống hệt nhau, làm sao chúng tách ra để thực sự bổ sung chi tiết
    thay vì vẫn cứ trùng khít mãi? \textbf{Trả lời:} ngay sau khi sinh ra, hai Gaussian đi qua \textbf{optimizer độc
    lập} (mỗi Gaussian có trạng thái Adam --- \texttt{exp\_avg}, \texttt{exp\_avg\_sq} --- riêng của nó). Ở những bước
    gradient descent tiếp theo, dù xuất phát từ cùng một điểm, hai Gaussian này chịu nhiễu số học và các tương tác nhỏ
    khác nhau với các Gaussian lân cận $\Rightarrow$ dần dần bị "đẩy" ra hai hướng khác nhau một cách tự nhiên, giúp phủ
    tốt hơn vùng đang thiếu chi tiết.
    \item Các bộ đếm thống kê (\texttt{xyz\_gradient\_accum}, \texttt{denom}, v.v.) của bản sao mới được \textbf{reset về
    $0$} ngay sau khi sinh --- lịch sử theo dõi của nó hoàn toàn tách biệt với Gaussian cha kể từ thời điểm này.
\end{itemize}
\begin{center}
\includegraphics[width=0.30\linewidth]{fig_08_clone_before_after.png}
\end{center}
\end{frame}
